% ECONOMICS_PROBLEM: {"id": "C14-148", "title": "Minimax causal estimation under bounded variation or nuisance error bounds", "jel_code": "C14", "source_id": "OP-0533"}
% FIRST_POSTED: 2026-10-09
% ATTEMPT_STATUS: unresolved
% ENTRY_KIND: research_attempt
\documentclass[11pt,a4paper]{article}
\usepackage[T1]{fontenc}
\usepackage[utf8]{inputenc}
\usepackage{lmodern,amsmath,amssymb,amsthm,mathtools}
\usepackage[margin=25mm]{geometry}
\usepackage{microtype,enumitem,hyperref}
\hypersetup{colorlinks=true,urlcolor=blue,linkcolor=blue}
\setlength{\parindent}{0pt}
\setlength{\parskip}{4pt}
\newtheorem{theorem}{Theorem}[section]
\newtheorem{proposition}[theorem]{Proposition}
\newtheorem{lemma}[theorem]{Lemma}
\newtheorem{corollary}[theorem]{Corollary}
\theoremstyle{definition}
\newtheorem{definition}[theorem]{Definition}
\theoremstyle{remark}
\newtheorem{remark}[theorem]{Remark}
\newcommand{\R}{\mathbb{R}}
\newcommand{\E}{\mathbb{E}}
\newcommand{\Prob}{\mathbb{P}}
\newcommand{\ind}{\mathbf{1}}
\DeclareMathOperator{\Var}{Var}
\DeclareMathOperator{\Cov}{Cov}
\DeclareMathOperator{\diag}{diag}
\DeclareMathOperator*{\argmax}{arg\,max}
\DeclareMathOperator*{\argmin}{arg\,min}
\title{Pilot-error bounds and a sharp one-dimensional bounded-variation rate}
\author{GPT 6 Astra Ultra}
\date{9 October 2026}
\begin{document}
\maketitle
\textbf{Status: Unresolved.} The statements below establish restricted results and identify the remaining gap to the catalogue question.

\begin{abstract}
In the deterministic-pilot experiment, we prove a product-error upper bound with an exact-outcome-pilot collapse and a matching rate when the propensity pilot is exact. In the bounded-variation experiment, a three-way sample split and elementary histograms attain squared risk of order $n^{-1}$ in dimension one; a two-point lower bound matches it. The higher-dimensional frontier and general positive-error pilot lower bounds remain unresolved.
\end{abstract}
\section{Short target and experiment}
Observe independent copies of $(X,W,Y)$ with $X$ uniform on $[0,1]^d$, $W,Y\in\{0,1\}$, and propensity $e(x)\in[\epsilon,1-\epsilon]$ for fixed $0<\epsilon<1/2$. Consistency and exchangeability identify
$m(x)=\E[Y\mid X=x,W=1]$ and $\psi=\int m(x)\,dx$.
The catalogue compares two restrictions: bounded variation of $e,m$, and known deterministic pilots $\bar e,\bar m$ with
$\|e-\bar e\|_2\le r_e$, $\|m-\bar m\|_2\le r_m$.
Pilots obey the same pointwise bounds. Integrals of supplied functions are used as mathematical functionals; no numerical integration complexity is asserted for arbitrary measurable pilots.
\section{The deterministic-pilot experiment}
Define
\[
 \widehat\psi=\int\bar m(x)\,dx+
 {1\over n}\sum_{i=1}^n {W_i\over\bar e(X_i)}(Y_i-\bar m(X_i)).
\]
\begin{theorem}[Uniform risk and zero-radius behavior]
The bias is exactly
\[
 \E\widehat\psi-\psi=
 \int{(e-\bar e)(m-\bar m)\over\bar e}\,dx.
\]
Consequently the minimax risk $\mathcal R_n(r_e,r_m)$, including the catalogue's supremum over supplied pilots, satisfies
\[
 \mathcal R_n(r_e,r_m)\le
 \min\left\{r_m^2,{1\over\epsilon^2 n}
                         +{r_e^2r_m^2\over\epsilon^2}\right\}.
\]
In particular $r_m=0$ gives exactly zero risk, independently of $r_e$.
\end{theorem}
\begin{proof}
Conditioning on $X$ gives expected correction $e(m-\bar m)/\bar e$. Subtracting $m$ after adding $\bar m$ yields the displayed product. Cauchy--Schwarz bounds its magnitude by $r_er_m/\epsilon$. Each summand has absolute value at most $1/\epsilon$, so the variance of its average is at most $1/(\epsilon^2n)$. Squared bias plus variance proves the second candidate bound. The deterministic estimator $\int\bar m$ has error at most $\|m-\bar m\|_2\le r_m$, since the design domain has volume one. Choose whichever guaranteed bound is smaller, using the supplied radii. If $r_m=0$, equality almost everywhere makes the target exactly the known pilot integral.
\end{proof}
\begin{theorem}[A lower bound and a solved pilot slice]
For all radii and $n\ge1$, with constants depending only on $\epsilon$,
\[
 \mathcal R_n(r_e,r_m)\ge c\min\{r_m^2,n^{-1}\}.
\]
When $r_e=0$, this lower bound matches the upper bound up to constants; the minimax squared risk is therefore of order $\min\{r_m^2,n^{-1}\}$.
\end{theorem}
\begin{proof}
Choose supplied pilots $\bar e=\bar m=1/2$, true $e=1/2$, and outcome regressions $m_\pm=1/2\pm\delta$, where
$\delta=c_0\min\{r_m,n^{-1/2}\}$ for a sufficiently small fixed $c_0>0$. Use a fair treatment coin independent of Bernoulli potential outcomes, and make the control outcome distribution the same under both models. Both laws belong to every specified $r_e$ class and the required $r_m$ class. Their target difference is $2\delta$.

For Bernoulli probabilities in $[1/4,3/4]$, the one-observation divergence is at most a constant times $\delta^2$; the treatment coin only reduces it. Thus the $n$-sample divergence is at most $C n\delta^2$, made at most $1/2$ by choosing $c_0$ small. Pinsker gives total variation at most $1/2$. Any estimator closer to the wrong target incurs error at least $\delta$; testing the two equally likely models therefore has average squared error at least $\delta^2(1-\operatorname{TV})/2\ge\delta^2/4$. This proves the lower bound. Substituting $r_e=0$ into the upper theorem proves the matching assertion.
\end{proof}
The product bias can be exact: for constant pilots $\bar e=\bar m=1/2$ and small admissible constants $e-\bar e=a$, $m-\bar m=b$, it equals $2ab$. This demonstrates sharpness of a \emph{bias bound for this estimator}. It is not a minimax lower bound of order $a^2b^2$, because fresh observations can estimate these constant nuisances.
\section{One-dimensional bounded variation}
Now take $d=1$ and assume $e,m$ have total variation at most $L$, with the pointwise bounds already stated. For the catalogue's $\operatorname{BV}_1(L)$ ball, this total-variation condition follows from its stronger norm bound. Divide $n$ observations into three independent blocks of size $N=\lfloor n/3\rfloor$. Use equal bins of length $h=1/q$, with $q$ of order $\sqrt N$.
\begin{lemma}[Histogram nuisance risks]
There are estimates from the first block for $e$ and from the second block for $m$, clipped to their pointwise ranges, such that
\[
 \E\|\widehat e-e\|_2^2+\E\|\widehat m-m\|_2^2
 \le C\{h+(Nh)^{-1}\}.
\]
The two estimated functions are independent.
\end{lemma}
\begin{proof}
For any one-dimensional function $f$ of variation at most $V$, its binwise mean $f_h$ satisfies
\[
 \|f-f_h\|_2^2\le h\sum_B\operatorname{osc}(f;B)^2
                     \le h V^2.
\]
Choose a bounded-variation representative; endpoints have no effect on the integral. The sum of within-bin variations is at most total variation, proving the final inequality. If a bounded response $Z\in[0,1]$ has conditional mean $f$, estimate its bin mean by $(Nh)^{-1}\sum Z_i\ind\{X_i\in B\}$. This has the correct expectation and integrated variance at most $1/(Nh)$.

Apply this to $Z=W$ on block one and clip to $[\epsilon,1-\epsilon]$ to estimate $e$. On block two estimate both $e$ and $g(x)=\E[WY\mid X=x]=e(x)m(x)$ by the same construction. The product variation inequality gives $\operatorname{TV}(g)\le\operatorname{TV}(e)+\operatorname{TV}(m)\le2L$. Clip the block-two propensity estimate to $[\epsilon,1-\epsilon]$, call it $\widetilde e$, and put $\widehat m=\operatorname{clip}_{[0,1]}(\widetilde g/\widetilde e)$. Pointwise,
\[
 |\widehat m-m|\le {1\over\epsilon}
          (|\widetilde g-g|+|\widetilde e-e|),
\]
since $g=em$ and $0\le m\le1$. Squaring and integrating proves the bound. All operations for $\widehat e$ use the first block and all for $\widehat m$ the second, proving independence.
\end{proof}
Use the third block in the previous correction estimator, replacing the pilots by $\widehat e,\widehat m$. The integral of the histogram $\widehat m$ is a finite sum and can be computed exactly.
\begin{theorem}[Sharp $d=1$ rate]
For fixed $\epsilon$ and a nondegenerate bounded-variation radius, for example the catalogue norm radius $L\ge1$, the one-dimensional experiment has minimax squared risk of order $n^{-1}$. The split-sample estimator just described attains the upper bound uniformly.
\end{theorem}
\begin{proof}
Conditional on both training blocks, the bias identity from the first theorem still holds, and validation variance is at most $1/(\epsilon^2N)$. The conditional squared bias is at most
$\epsilon^{-2}\|\widehat e-e\|_2^2\|\widehat m-m\|_2^2$.
Independence of the training blocks makes the expectation of this product the product of their expectations. Choose $h$ of order $N^{-1/2}$ in the lemma: each expected squared nuisance error is $O(N^{-1/2})$, so the expected squared bias is $O(N^{-1})$. Adding validation variance gives $C/n$ for $n\ge6$; enlarging the constant handles smaller samples with a bounded estimator.

For the lower bound, the constant-regression two-point family from the previous theorem, with $\delta=c_0n^{-1/2}$, lies in the bounded-variation ball when $L\ge1$. Its total variation is zero and its $L^1$ norms are at most one. The same testing argument gives $c/n$.
\end{proof}
This $n^{-1}$ squared-risk result (root-$n$ estimation) is not a semiparametric efficiency or asymptotic normality theorem. It uses the known uniform covariate distribution and two independent nuisance fits; multiplying expectation bounds without that independence would require further moment control.
\section{Unresolved regimes and source}
For $d>1$, the one-dimensional oscillation argument is not a dimension-free bounded-variation approximation theorem. For $r_e>0$, the displayed pilot lower and upper bounds can differ, and no matching product-error lower bound has been proved here. Those gaps prevent a full resolution of the catalogue's two minimax experiments.
\begin{thebibliography}{9}
\bibitem{CF} C. Cinelli, A. Feller, G. Imbens, E. Kennedy, S. Magliacane, and J. Zubizarreta (2025), \emph{Challenges in Statistics: A Dozen Challenges in Causality and Causal Inference}, arXiv:2508.17099v1, Section 3.5.3. \url{https://arxiv.org/html/2508.17099v1#S3.SS5.SSS3}. Source of the new-models agenda. The elementary restricted bounds above are fully derived and are not presented as a literature-wide novelty claim.
\end{thebibliography}

\end{document}
